\subsection{唯一性定理}

命题 4. 设 \((\mathfrak{g},\mathfrak{h},\mathrm{B},(X_{\alpha})_{\alpha\in\mathrm{B}})\) 是一个带标架的半单李代数。设 \((n(\alpha,\beta))_{\alpha,\beta\in\mathrm{B}}\) 与 \((X_{\alpha})_{\alpha\in\mathrm{B}\cup(-\mathrm{B})}\) 是相应的 Cartan 矩阵与生成族。

\begin{enumerate}
\def\labelenumi{(\roman{enumi})}
\item
  族 \(((X_{\alpha}, H_{\alpha}, X_{-\alpha}))_{\alpha \in \mathrm{B}}\) 连同第 3 节的关系式 (16) 至 (22) 构成 \(\mathfrak{g}\) 的一个呈现。
\item
  族 \((X_{\alpha})_{\alpha \in \mathrm{B}}\) 连同第 3 节的关系式 (21) 构成 \(\mathfrak{g}\) 的由 \((X_{\alpha})_{\alpha \in \mathrm{B}}\) 所生成的子代数的一个呈现。
\end{enumerate}

设 R 是 \((\mathfrak{g},\mathfrak{h})\) 的根系。把第 2 节与第 3 节中的构造应用于 R 与 B ，我们得到一些对象，将用 \(a', g', X'_{\alpha}, H'_{\alpha}, \ldots\) 来记它们，以代替 \(a, g, X_{\alpha}, H_{\alpha}, \ldots\)

存在同态 \(\varphi\) （从李代数 \(\mathfrak{g}'\) 到李代数 \(\mathfrak{g}\) ），使 \(\varphi(X_{\alpha}') = X_{\alpha}\) ，\(\varphi(H_{\alpha}') = H_{\alpha}\) ，\(\varphi(X_{-\alpha}') = X_{-\alpha}\) 对一切 \(\alpha \in B\) 成立（命题 1）。由于 \(\dim \mathfrak{g}' = \operatorname{Card} R + \operatorname{Card} B = \dim \mathfrak{g}\) ，\(\varphi\) 是双射。这就证明了 (i)。

\(\mathfrak{g}' = \mathfrak{a}' / (\mathfrak{n}'\oplus \theta'\mathfrak{n}') = (\mathfrak{a}_+^\prime \oplus \mathfrak{a}^{\prime 0}\oplus \mathfrak{a}_-^\prime) / (\mathfrak{n}'\oplus \theta'\mathfrak{n}')\) 中由 \((X_{\alpha}^{\prime})_{\alpha \in \mathrm{B}}\) 生成的子代数可与 \(\mathfrak{a}_+^\prime /\mathfrak{n}'\) 等同。鉴于命题 3 与 \(\mathfrak{n}'\) 的定义，这就证明了 (ii)。

推论. 每个带标架的半单李代数都可由一个带标架的半单 \(\mathbf{Q}\) -李代数经由基域从 \(\mathbf{Q}\) 到 \(k\) 的扩张而得到。

这由该命题立即得出。

定理 2. 设 \((\mathfrak{g},\mathfrak{h},\mathrm{B},(X_{\alpha})_{\alpha \in \mathrm{B}})\) 与 \((\mathfrak{g}',\mathfrak{h}',\mathrm{B}',(X'_{\alpha})_{\alpha \in \mathrm{B}'})\) 是带标架的半单李代数，设 R 与 R' 是 \((\mathfrak{g},\mathfrak{h})\) 与 \((\mathfrak{g}',\mathfrak{h}')\) 的根系，设 \((n(\alpha ,\beta))_{\alpha ,\beta \in \mathrm{B}}\) （相应地 \((n^{\prime}(\alpha ,\beta))_{\alpha ,\beta \in \mathrm{B}'}\) ）是 R （相应地 R' ）相对于 B （相应地 B' ）的 Cartan 矩阵，并设 \(\Delta\) （相应地 \(\Delta'\) ）是 R （相应地 R' ）相对于 B （相应地 B' ）的邓金图。

\begin{enumerate}
\def\labelenumi{(\roman{enumi})}
\item
  若 \(\varphi\) 是从 \(\mathfrak{h}^*\) 到 \(\mathfrak{h}'^*\) 的同构，使 \(\varphi(\mathrm{R}) = \mathrm{R}'\) 且 \(\varphi(\mathrm{B}) = \mathrm{B}'\) ，则存在唯一的同构 \(\psi\) （从 \((\mathfrak{g}, \mathfrak{h}, \mathrm{B}, (X_{\alpha})_{\alpha \in \mathrm{B}})\) 到 \((\mathfrak{g}', \mathfrak{h}', \mathrm{B}', (X_{\alpha}')_{\alpha \in \mathrm{B}'})\) ），使 \(\psi|_{\mathfrak{h}} = {}^t\varphi^{-1}\) 。
\item
  若 \(f\) 是从 \(\mathrm{B}\) 到 \(\mathrm{B}'\) 的双射，使 \(n'(f(\alpha), f(\beta)) = n(\alpha, \beta)\) 对一切 \(\alpha, \beta \in \mathrm{B}\) 成立，则存在从 \((\mathfrak{g}, \mathfrak{h}, \mathrm{B}, (X_{\alpha})_{\alpha \in \mathrm{B}})\) 到 \((\mathfrak{g}', \mathfrak{h}', \mathrm{B}', (X_{\alpha}')_{\alpha \in \mathrm{B}'}).\) 的同构
\item
  若存在从 \(\Delta\) 到 \(\Delta'\) 的同构，则存在从 \((\mathfrak{g},\mathfrak{h},\mathrm{B},(X_{\alpha})_{\alpha \in \mathrm{B}})\) 到 \((\mathfrak{g}',\mathfrak{h}',\mathrm{B}',(X_{\alpha}')_{\alpha \in \mathrm{B}'})\) 的同构。
\end{enumerate}

这由命题 4 (i) 立即得出（对部分 (iii) 要用到第六章 §4 第 2 节 命题 1）。

概要. 对每个可分裂半单李代数 g 都联系着一个邓金图，它在同构意义下确定 g （定理 2 (iii)）。此图非空且连通，当且仅当 g 是单的（§3 第 2 节 命题 6 的推论 1）。由第 3 节的定理 1 以及第六章 §4 第 2 节 定理 3，可分裂单李代数就是 \(A_{l}\) 型（ \(l \geq 1\) ）、\(B_{l}\) 型（ \(l \geq 2\) ）、\(C_{l}\) 型（ \(l \geq 3\) ）、\(D_{l}\) 型（ \(l \geq 4\) ）、\(E_{6}\) 型、\(E_{7}\) 型、\(E_{8}\) 型、\(F_{4}\) 型、\(G_{2}\) 型的李代数。此列表中任何两个代数都互不同构。

命题 5. 设 \((\mathfrak{g},\mathfrak{h},\mathrm{B},(X_{\alpha})_{\alpha \in \mathrm{B}})\) 是一个带标架的半单李代数，\((X_{\alpha})_{\alpha \in \mathrm{B}\cup (-\mathrm{B})}\) 是相应的生成族。存在唯一的自同构 \(\theta\) （\(\mathfrak{g}\) 的），使 \(\theta (X_{\alpha}) = X_{-\alpha}\) 对一切 \(\alpha \in \mathrm{B}\cup (-\mathrm{B})\) 成立。我们有 \(\theta^2 = \mathrm{Id}_{\mathfrak{g}}\) ，且 \(\theta (h) = -h\) 对一切 \(h\in \mathfrak{h}\) 成立。

唯一性是显然的，因为 \((X_{\alpha})_{\alpha\in\mathrm{B}\cup(-\mathrm{B})}\) 生成李代数 g。鉴于命题 4，\(\theta\) 的存在性由第 3 节中定理 1 之前所述的内容得出。

推论. 设 \((\mathfrak{g},\mathfrak{h})\) 是一个分裂半单李代数。则 \((\mathfrak{g},\mathfrak{h})\) 拥有一个谢瓦莱系（§2 第 4 节 定义 3）。

设 R 是 \((\mathfrak{g},\mathfrak{h})\) 的根系。对一切 \(\alpha\in R\) ，设 \(X_{\alpha}\) 是 \(g^{\alpha}\) 的一个非零元素。假设诸 \(X_{\alpha}\) 的选取使 \([X_{\alpha},X_{-\alpha}]=-H_{\alpha}\) 对一切 \(\alpha\in R\) 成立（§2 第 4 节 引理 2）。设 B 是 R 的一个基，\(\theta\) 是 g 的自同构，使 \(\theta(X_{\alpha})=X_{-\alpha}\) 对一切 \(\alpha\in\mathrm{B}\cup(-\mathrm{B})\) 成立。我们有 \(\theta|_{\mathfrak{h}}=-\mathrm{Id}_{\mathfrak{h}}\) 。于是，对一切 \(\alpha\in R\) ，存在 \(t_{\alpha}\in k^{*}\) 使 \(\theta X_{\alpha}=t_{\alpha}X_{-\alpha}\) 。我们有

\[
\begin{array}{c} t _ {\alpha} t _ {- \alpha} H _ {\alpha} = [ t _ {\alpha} X _ {- \alpha}, t _ {- \alpha} X _ {\alpha} ] = [ \theta X _ {\alpha}, \theta X _ {- \alpha} ] = \theta ([ X _ {\alpha}, X _ {- \alpha} ]) \\ = \theta (- H _ {\alpha}) = H _ {\alpha} \end{array}
\]

故 \(t_{\alpha}t_{-\alpha} = 1\) 对一切 \(\alpha \in \mathrm{R}\) 成立。按 §2 第 4 节引入诸 \(\mathrm{N}_{\alpha \beta}\) 。若 \(\alpha, \beta, \alpha + \beta \in \mathrm{R}\) ，

\[
\begin{array}{r} \mathrm{N} _ {- \alpha , - \beta} t _ {\alpha} t _ {\beta} X _ {- \alpha - \beta} = t _ {\alpha} t _ {\beta} [ X _ {- \alpha}, X _ {- \beta} ] = [ \theta X _ {\alpha}, \theta X _ {\beta} ] = \theta ([ X _ {\alpha}, X _ {\beta} ]) \\ = \mathrm{N} _ {\alpha \beta} \theta X _ {\alpha + \beta} = \mathrm{N} _ {\alpha \beta} t _ {\alpha + \beta} X _ {- \alpha - \beta} \end{array}
\]

于是，鉴于 §2 第 4 节 引理 4，

\[
(q + 1) ^ {2} t _ {\alpha} t _ {\beta} = \mathrm{N} _ {\alpha \beta} ^ {2} t _ {\alpha + \beta}
\]

其中 \(q\) 是一个整数。由此推出，若 \(t_{\alpha}\) 与 \(t_{\beta}\) 是 \(k^*\) 中的平方元，则 \(t_{\alpha + \beta}\) 也是。由于 \(t_{\alpha} = 1\) 对一切 \(\alpha \in \mathrm{B}\) 成立，第六章 §1 第 6 节 命题 19 证明 \(t_{\alpha}\) 对一切 \(\alpha \in \mathrm{R}\) 都是平方元。对一切 \(\alpha \in \mathrm{R}\) ，选取 \(u_{\alpha} \in k\) 使 \(u_{\alpha}^{2} = t_{\alpha}\) 。这一选取可以做到使 \(u_{\alpha}u_{-\alpha} = 1\) 对一切 \(\alpha \in \mathrm{R}\) 成立。置 \(X_{\alpha}' = u_{\alpha}^{-1}X_{\alpha}\) 。则对一切 \(\alpha \in \mathrm{R}\) ，

\[
X _ {\alpha} ^ {\prime} \in \mathfrak {g} ^ {\alpha}, [ X _ {\alpha} ^ {\prime}, X _ {- \alpha} ^ {\prime} ] = [ X _ {\alpha}, X _ {- \alpha} ] = - H _ {\alpha},
\]

且 \(\theta(X_{\alpha}^{\prime}) = \theta(u_{\alpha}^{-1}X_{\alpha}) = u_{\alpha}^{-1}t_{\alpha}X_{-\alpha} = u_{\alpha}X_{-\alpha} = u_{\alpha}u_{-\alpha}X_{-\alpha}^{\prime} = X_{-\alpha}^{\prime}\) ，于是 \((X_{\alpha}^{\prime})_{\alpha \in \mathrm{R}}\) 是 \((\mathfrak{g},\mathfrak{h})\) 的一个谢瓦莱系。

